%This document uses PGF 3.0 graph package.
%Compile it with lualatex.

\documentclass[a4paper]{exam}

\usepackage[utf8]{inputenc}
\usepackage{fontspec}
\usepackage[portuguese]{babel}
\usepackage{amsmath,amssymb, alltt}
\usepackage[margin=1.5cm,nohead,nofoot]{geometry}
\usepackage{multicol}
\usepackage{mathtools}
\usepackage{tikz}
\usepackage{graphicx}
\usepackage{caption}
\usepackage{subcaption}
\usepackage{float}
\usepackage{color}


\usepackage[ruled,vlined,noend,linesnumbered]{algorithm2e}
\setlength{\algoheightrule}{0pt}
\setlength{\algotitleheightrule}{0pt}




\DeclarePairedDelimiter{\ceil}{\lceil}{\rceil}
\DeclarePairedDelimiter{\floor}{\lfloor}{\rfloor}

\footer{}{}{}

\checkboxchar{$\Box$}
\checkedchar{$\blacksquare$}
\bracketedpoints
\shadedsolutions
\renewcommand{\solutiontitle}{\noindent\textbf{Solution:}\par\noindent}
\definecolor{SolutionColor}{gray}{0.8}
%\def\explanation#1{%
\def\explanation#1{%
  \ifprintanswers\mdseries\hspace{1em}\colorbox{SolutionColor}‌​{\strut #1}
  \ifprintanswers\mdseries\hspace{1em}\colorbox{SolutionColor}{\strut #1}
  \else 
  \enspace\hrulefill
  \fi}

\renewcommand{\solutiontitle}{\noindent\textbf{Solução:}\par\noindent}


\begin{document}

%Descomente a linha a seguir para imprimir as soluções.
\printanswers

\begin{center}
\begin{minipage}[c][1.5cm][c]{1.5cm}
\includegraphics[height=1.5cm]{logo.jpg}
\end{minipage}
\begin{minipage}[c][1.5cm][c]{12cm} 
\textsc{\Large Universidade Estadual de Maringá} \\
Departamento de Informática -- Projeto e Análise de Algoritmos \\
Prof. Daniel Kikuti
\end{minipage}
\end{center}

\begin{center}
\section*{Trabalho 1}
\end{center}

\bracketedpoints
\pointname{}
\pointformat{[Valor: \thepoints]}

\begin{questions}
  %\large
  %01
  \question Para cada par de funções $f(n)$ e $g(n)$ na tabela a
  seguir, indique se $f(n)$ pertence a $O(g(n))$, $\Omega(g(n))$
  ou $\Theta(g(n))$. Considere que $k \geq 1$, $\epsilon > 0$ e
  $c > 1$. Justifique sua resposta.
  \begin{center}
    \begin{tabular}{|c|c|c|c|c|c|}
      \hline
      & $f(n)$ & $g(n)$ & $f(n) = O(g(n))$? & $f(n) = \Omega(g(n))$? & $f(n) = \Theta(g(n))$? \\
      \hline
      a) & $\lg^k n$  &  $n^{\epsilon}$   & & X & \\ \hline
      b) & $n^k$      &  $c^n$          & X & & \\ \hline
      c) & $2^n$      &  $2^{n/2}$       & & X & \\ \hline
      d) & $n^{\lg c}$ &  $c^{\lg n}$      & & & X \\ \hline
    \end{tabular}
  \end{center}
  \begin{solution}
    a) $\displaystyle\lim_{n\to\infty} \frac{\lg^k n}{n^{\epsilon}} = \infty \to f(n) = \Omega(g(n))$ \\
    b) $\displaystyle\lim_{n\to\infty} \frac{n^k}{c^n} = 0 \to f(n) = O(g(n))$ \\
    c) $\displaystyle\lim_{n\to\infty} \frac{2^n}{2^{n/2}} = \infty \to f(n) = \Omega(g(n))$ \\
    d) $\displaystyle\lim_{n\to\infty} \frac{n^{\lg c}}{c^{\lg n}} = 1 \to f(n) = \Theta(g(n))$
  \end{solution}
  %02
  \question Para cada item a seguir, assinale \textbf{V}erdadeiro 
  ou \textbf{F}also. \textbf{Justifique} sua resposta usando as 
  definições de notação assintótica.
  \begin{parts}
    \part 
    \mbox{\begin{oneparcheckboxes}
      \correctchoice V \choice F
    \end{oneparcheckboxes}}
    Se $f(n) = \log_{16} n$ então $f(n) = \Theta(\lg n)$?
    \part 
    \mbox{\begin{oneparcheckboxes}
      \choice V \correctchoice F
    \end{oneparcheckboxes}}
    $2^{n + a} = \Theta(2^{2n})$? Onde $a \in \mathbb{N}$ é uma constante.
    \part
    \mbox{\begin{oneparcheckboxes}
      \correctchoice V \choice F
    \end{oneparcheckboxes}}
    $\frac{n^{2}}{4} - 3n - 16 = \Omega(n^2)$?
    \part
    \mbox{\begin{oneparcheckboxes}
      \correctchoice V \choice F
    \end{oneparcheckboxes}}
    $7n^2 + 13n = O(n^2)$.
  \end{parts}
  \begin{solution}
    a) $\displaystyle\lim_{n\to\infty} \frac{\log_{16} n}{\lg n} = \frac{1}{4} \to f(n) = \Theta(g(n))$ \\
    b) $\displaystyle\lim_{n\to\infty} \frac{2^{n+a}}{2^{2n}} = 0 \to f(n) = 2^{n+a} \in O(2^{2n})$ \\
    c) $\displaystyle\lim_{n\to\infty} \frac{\frac{n^2}{4} -3n - 16}{n^2} = \frac{1}{4} \to f(n) = \frac{n^2}{4} -3n - 16 
        \in \Theta(n^2) =  O(n^2) \land  \Omega(n^2)$ \\
    d) $\displaystyle\lim_{n\to\infty} \frac{7n^2 +13n}{n^2} = 7 \to f(n) = \frac{7n^2 +13n}{n^2}
        \in \Theta(n^2) =  O(n^2) \land  \Omega(n^2)$
  \end{solution}
  %03
  \question Mostre usando as definições de notação assintótica:
  \begin{parts}
    \part $\frac{n}{2}\lg(\frac{n}{2}) = \Omega(n \lg n)$.
    \part $n^2 + 1000n = O(n^2)$.
    \part $2^{n+1} = \Theta(2^{n})$.
  \end{parts}
  \begin{solution}
    Escreva a solução aqui.
  \end{solution}
  %04
  \question Expresse as seguintes funções em termos da notação $\Theta$.
  \begin{parts}
    \part $2n + 3 \log ^{100} n$.
    \part $7n^3 + 1000n \log n + 3n$.
    \part $3n^{1.5} + (\sqrt{n})^3 \log n$.
    \part $2^n + 100^n + n!$.
  \end{parts}
  \begin{solution}
    Escreva a solução aqui.
  \end{solution}
  %05
  \question É comum usar a notação $f(n) \prec g(n)$ para denotar que
  $f(n) \in o(g(n))$. Use esta notação para expressar a
  hierarquia de classes de complexidade das seguintes funções:
  $\sqrt{n}$, $2^{n^2}$, $n$, $\lg n$, $1$, $\lg \lg n$, $n!$, $n^2$, $n^{3/4}$, $2^n$, $n \lg n$.
  \begin{solution}
    Escreva a solução aqui.
  \end{solution}
  
  
  %06
  \question Sejam $f(n)$ e $g(n)$ funções positivas. Informe se a
  afirmação é verdadeira ou falsa e justifique.
  \begin{parts}
    \part $f(n) = O(g(n))$ implica $g(n) = O(f(n))$.
    \part $f(n) + g(n) = \Theta(min(f(n),g(n)))$
    \part $f(n) = O(g(n))$ implica $g(n) = \Omega(f(n))$
    \part $f(n) = \Theta(f(n/2))$
  \end{parts}
  \begin{solution}
    Escreva a solução aqui.
  \end{solution}
  %07
  \question Mostre uma função $f(n)$ tal que $f(n) \not\in \Omega(f(n+1))$.
  \begin{solution}
    Escreva a solução aqui.
  \end{solution}
  %08
  \question Mostre que $\sum_{i=1}^{n} \lg i = \Theta(n \lg n)$.
  \begin{solution}
    Escreva a solução aqui.
  \end{solution}
  %09
  \question Mostre que $n! = O(2^{n^2})$.
  \begin{solution}
    Escreva a solução aqui.
  \end{solution}
  %10
  \question Seja $p(n) = \sum_{i=0}^{k} a_{i}n^{i}$
  (polinômio em $n$ de grau $k$), onde $k$ é um inteiro não-negativo, $a_i$
  é uma constante e $a_k > 0$, mostre que $p(n) = \Theta(n^k)$.
  \begin{solution}
    Usando limite, devemos mostrar que: L = $\displaystyle\lim_{n\to\infty} \frac{\sum_{i=0}^{k} a_{i}n^{i}}{n^k}, 0 < L < \infty$\\

    $ L = \displaystyle\lim_{n\to\infty} \frac{\sum_{i=0}^{k} a_{i}n^{i}}{n^k} \\
        = \displaystyle\lim_{n\to\infty} \frac{a_0 n^0 + a_1 n^1 + a_2 n^2 + ... + a_k n^k}{n^k} \\
        = \displaystyle\lim_{n\to\infty} \frac{a_0 n^0}{n^k} + 
          \displaystyle\lim_{n\to\infty} \frac{a_1 n^1}{n^k} + 
          \displaystyle\lim_{n\to\infty} \frac{a_2 n^2}{n^k} + ... + 
          \displaystyle\lim_{n\to\infty} \frac{a_k n^k}{n^k} \\
        = 0 + 0 + 0 + a_k \\
        = a_k.$ \\
    Como $a_k$ > 0 => $p(n) = \Theta(n^k)$.
  \end{solution}
  %11
  \question Papai Noel resolveu antecipar seu presente de
  Natal. Sabendo que você foi um(a) bom(a) menino(a), ele escreveu um
  algoritmo para que você analise e ganhe uns pontinhos na prova de
  PAA. Sua tarefa é simples. Dado um inteiro $n$ como entrada,
  expresse por meio de notação assintótica a quantidade de ``Ho!''s
  que será impressa pelo algoritmo (use a notação $\Theta$).
  
  \begin{algorithm}[H]
    \NoCaptionOfAlgo
    \DontPrintSemicolon

    $i \leftarrow 1$\;
    \While{$i \leq n$} {
      \For {$j \leftarrow i$ \KwTo $2i - 1$}{
        print ``Ho!'';
      }
      $i \leftarrow 2i$\;
    }
    \caption{\textsc{Feliz-Natal}($n$)}
  \end{algorithm}
  \begin{solution}
    Escreva a solução aqui.
  \end{solution}
  %12
  \question Dado um inteiro $n$ (assuma que $n=2^k$, tal que
  $k$ é um número inteiro positivo) e $expr$ (que corresponde a uma
  expressão a ser impressa), informe, por meio de notação assintótica,
  a quantidade de mensagens que o algoritmo a seguir irá
  imprimir. Dê sua resposta em função de $n$.

  \begin{algorithm}[H]
    \NoCaptionOfAlgo
    \DontPrintSemicolon

    \While {$n \geq 1$}{
      \For {$j \gets 1$ \KwTo $n$}{
        print $expr$\;
      }
      $n \gets n / 2$
    }
    \caption{\textsc{Prog}($n$, $expr$)}
  \end{algorithm}
  \begin{solution}
    Escreva a solução aqui.
  \end{solution}
  %13
  \question Seja $count$ o número total de iterações feitas pelo
  algoritmo a seguir para uma entrada $n$. Informe, usando
  notação assintótica, o valor de $count$ em função de $n$.

  \begin{algorithm}[H]
    \NoCaptionOfAlgo
    \DontPrintSemicolon

    $count \gets 0$\;
    \For {$i \gets 1$ \KwTo $n$}{
      \For {$j \gets 1$ \KwTo $\floor{n/i}$}{
        $count \gets count + 1$\;
      }
    }
    \Return $count$\;
    \caption{\textsc{Count}($n$)}
  \end{algorithm}
  \begin{solution}
    Escreva a solução aqui.
  \end{solution}
  %14
  \question Seja $count$ o número total de iterações feitas pelo
  algoritmo a seguir para uma entrada $n$ (considere que $n = 2^{2^{k}}$,
  para algum inteiro positivo $k$). Informe, usando
  notação assintótica, o valor de $count$ em função de $n$.
  
  \begin{algorithm}[H]
    \NoCaptionOfAlgo
    \DontPrintSemicolon

    $count \gets 0$\;
    \For {$i \gets 1$ \KwTo $n$}{
      \For {$j \gets 2; j \leq n; j \gets j^2$}{
        $count \gets count + 1$\;
      }
    }
    \Return $count$\;
    \caption{\textsc{Count}($n$)}
  \end{algorithm}
  \begin{solution}
    Escreva a solução aqui.
  \end{solution}
  %15
  \question Dado um inteiro $n \geq 0$ como entrada, muitos
  afirmam que o algoritmo a seguir é capaz de medir o desespero na
  prova de PAA. Outros afirmam que o algoritmo mede a alegria. Para
  o professor, não interessa o que o algoritmo mede. O objetivo
  desta questão é avaliar se o aluno é capaz de encontrar uma
  fórmula fechada que representa o valor final de $x$ em função do
  valor de entrada $n$. Em outras palavras, uma função que
  representa quantas vezes a linha \ref{incrementa} será executada.

  \begin{algorithm}[H]
    \NoCaptionOfAlgo
    \DontPrintSemicolon

    $x \gets 0$\;
    \For {$i \gets 1$ \KwTo $n$}{
      \For {$j \gets i + 1$ \KwTo $n$}{
        \For {$k \gets 1$ \KwTo $j - i$}{
          $x \gets x + 1$ \label{incrementa}\;
        }
      }
    }
    \caption{\textsc{Prog}($n$)}
  \end{algorithm}
  \begin{solution}
    Escreva a solução aqui.
  \end{solution}
  %16
  \question Resolva as seguintes recorrências (use o método da
  substituição):
  \begin{parts}
    \part $T(n) = 8T(n/2) + \Theta(n^2)$
    \part $T(n) = 7T(n/2) + \Theta(n^2)$
    \part $T(n) = T(n/4) + 1$
    \part $T(n) = 2T(n/2) + n \lg n$
  \end{parts}
  \begin{solution}
    Escreva a solução aqui.
  \end{solution}
  %17
  \question Use árvore de recorrência para estimar um limite superior
  para as seguintes recorrências. Assuma que $T(n)$ é uma constante
  para $n \leq 2$. Depois comprove usando o método de
  substituição.
  \begin{parts}
    \part $T(n/2) + T(n/4) + T(n/8) +n$
    \part $2T(n/4) + \sqrt{n}$
  \end{parts}
  \begin{solution}
    Escreva a solução aqui.
  \end{solution}
  %18
  \question Utilize o método de árvore de recursão para supor um
  limite assintótico superior para a recorrência \linebreak
  $T(n) = 3T(n-1) + 1$. Depois verifique pelo método de
  substituição que este limite está correto.
  \begin{solution}
    Escreva a solução aqui.
  \end{solution}
  %19
  \question Utilize o método de árvore de recursão para supor um
  limite assintótico superior para a recorrência \linebreak
  $T(n) = 2T(n/2) + n \lg n$. Depois verifique pelo método de
  substituição que este limite está correto. 
  \begin{solution}
    Escreva a solução aqui.
  \end{solution}
  %20
  \question Utilize o método de árvore de recursão para supor um
  limite assintótico superior para a recorrência \linebreak $T(n) = T(n/3) +
  T(2n/3) + \Theta(n)$. Depois verifique pelo método de
  substituição que este limite está correto.
  \begin{solution}
    Escreva a solução aqui.
  \end{solution}
  %21
  \question A recorrência $T(n) = 7T(n/2) + n^2$ descreve o
  tempo de execução de um algoritmo $A$. Um algoritmo alternativo $A'$ 
  tem um tempo de execução $T'(n) = aT'(n/4) + n^2$. Qual é o maior
  inteiro $a$ que faz com que $A'$ seja assintoticamente mais rápido
  que $A$?
  \begin{solution}
    Escreva a solução aqui.
  \end{solution}
  %22
  \question Use o método mestre para resolver as seguintes
  recorrências: 
  \begin{parts}
    \part $T(n) = 3T(n/2) + n \lg n$ 
    \part $T(n) = 3T(n/2) + n^2$
    \part $T(n) = 4T(n/2) + n^2$
    \part $T(n) = 4T(n/2) + n^2\sqrt{n}$
    \part $T(n) = 5T(n/5) + n$  
    \part $T(n) = 6T(n/3) + n^2$
    \part $T(n) = 9T(n/2) + n^3$
  \end{parts}
  \begin{solution}
    Escreva a solução aqui.
  \end{solution}
  %23
  \question Dadas as recorrências dos algoritmos $A$ e
  $B$, determine a complexidade de cada um deles e compare-os
  (informe se $A$ é assintoticamente mais rápido que $B$, $B$
  é assintoticamente mais rápido que $A$, ou ambos possuem a
  mesma complexidade assintótica).
  \begin{itemize}
  \item $T_A(n) = 27T_A(n/3) + n$
  \item $T_B(n) = 4T_B(n/2) + n^3$
  \end{itemize}
  \begin{solution}
    Escreva a solução aqui.
  \end{solution}
  %24
  \question Os três algoritmos a seguir resolvem um
  problema de tamanho $n$ por meio da técnica de divisão e conquista.
  Analise a complexidade de cada um deles e informe qual algoritmo é
  assintoticamente mais eficiente.
  \begin{itemize}
  \item Algoritmo $A$ resolve problemas dividindo-os em cinco
    subproblemas de tamanho $n/2$, recursivamente resolve cada
    subproblema e combina suas soluções em tempo linear para obter
    uma solução do problema original.
    \item Algoritmo $B$ resolve problemas dividindo-os em dois
      subproblemas de tamanho $n-1$, recursivamente resolve cada um
      dos subproblemas e combina as soluções em tempo constante para
      obter a solução do problema original.
    \item Algoritmo $C$ resolve problemas dividindo-os em nove 
      subproblemas de tamanho $n/3$, recursivamente resolve cada
      subproblema e combina suas soluções em tempo $O(n^2)$ para obter
      uma solução do problema original.
  \end{itemize}
  \begin{solution}
    Escreva a solução aqui.
  \end{solution}
  %25
  \question Os três algoritmos a seguir computam corretamente
  $x^n$ para $x>0$ e $n \geq 0$. Mostre que os três algoritmos estão
  corretos e analise a complexidade assintótica
  de cada um deles (em função de $n$) e informe qual deles é mais
  eficiente.
  \setlength{\algoheightrule}{0pt}
  \setlength{\algotitleheightrule}{0pt}
  \begin{center}
    \vspace{-1cm}
    \begin{minipage}[t][][c]{.25\textwidth}
      \begin{algorithm}[H]
        \NoCaptionOfAlgo
        \DontPrintSemicolon
        $resp \gets 1$\;
        $i \gets 0$\;
        \While {$i < n$} {
          $resp \gets resp * x$\;
          $i \gets i + 1$\;
        }
        \Return $resp$\;
        \caption{\textsc{Power1}($x, n$)}
      \end{algorithm}
    \end{minipage}%
  \begin{minipage}[t][][c]{.35\textwidth} 
      \begin{algorithm}[H]
        \NoCaptionOfAlgo
        \DontPrintSemicolon
        \If {$n = 0$}{
          \Return $1$\;
        }
        \Else {
            \Return \textsc{Power2}($x, n-1$) $*$ $x$\;
        }
        \caption{\textsc{Power2}($x,n$)}
      \end{algorithm}
    \end{minipage}%
    \begin{minipage}[t][][c]{.38\textwidth} 
      \begin{algorithm}[H]
        \NoCaptionOfAlgo
        \DontPrintSemicolon
        \If {$n = 0$}{
          \Return $1$\;
        }
        \ElseIf {($n$ mod $2$) = 0} {
          $aux \leftarrow $ \textsc{Power3}($x, n/2$)\;
          \Return $aux*aux$\;
        }
        \Else {
          \Return \textsc{Power3}($x, n-1$) $*$ $x$\;
        }
        \caption{\textsc{Power3}($x,n$)}
      \end{algorithm}
    \end{minipage}
  \end{center}
  \begin{solution}
    Escreva a solução aqui.
  \end{solution}
  %26
  \question Considere o algoritmo \textsc{Heapsort} descrito a
  seguir. Mostre que o algoritmo está correto usando o seguinte
  invariante de laço: ``No começo de cada iteração do laço
  \textbf{for} das linhas 2--5, o subvetor $A[1 \dots i]$ contém os $i$
  menores elementos de $A[1 \dots n]$, e o subvetor $A[i{+}1 \dots n]$ contém
  os $n-i$ maiores elementos de $A[1 \dots n]$ em ordem.

  \setlength{\algoheightrule}{0pt}
  \setlength{\algotitleheightrule}{0pt}
  
  \IncMargin{1em}
  \begin{algorithm}[H]
    \NoCaptionOfAlgo
    \DontPrintSemicolon

    \textsc{Build-Max-Heap}($A$)\;
    \For {$i \leftarrow A.length$ \KwTo $2$}{
      \textsc{swap}($A[1], A[i]$)\;
      $A.heap$-$size \leftarrow A.heap$-$size - 1$\;
      \textsc{Max-Heapify}($A, 1$)\;
    }

    \caption{\textsc{Heapsort}($A$)}
  \end{algorithm}
  \DecMargin{1em}
  \begin{solution}
    Escreva a solução aqui.
  \end{solution}
  %27
  \question Analise a complexidade dos seguintes algoritmos:
  \begin{parts}
    \part
    \begin{algorithm}[H]
      \NoCaptionOfAlgo
      \DontPrintSemicolon
      \If {$n = 1$}{
        \Return 1\;
      }
      \Else {
        \Return \textsc{f}($n - 1$) + \textsc{f}($n - 1$)\;
      }
      \caption{\textsc{f}($n$)}
    \end{algorithm}
    \begin{solution}
      Escreva a solução aqui.
    \end{solution}
    
    \part
    \begin{algorithm}[H]
      \NoCaptionOfAlgo
      \DontPrintSemicolon
      \If {$max < min$}{
        \Return -1\;
      }
      $mid \leftarrow min + ((max - min)/2)$\;
      \If {$A[mid] > key$}{
        \Return \textsc{Busca}($A, key, min, mid-1$)\;
      }
      \ElseIf {A[mid] < key} {
        \Return \textsc{Busca}($A, key, mid+1, max$)\;
      }
      \Else {
        \Return $mid$\;
      }
      \caption{\textsc{Busca}($A[], key, min, max$)}
    \end{algorithm}
    \begin{solution}
      Escreva a solução aqui.
    \end{solution}
  
    \part
    \begin{algorithm}[H]
      \NoCaptionOfAlgo
      \DontPrintSemicolon
      \If {$n > 1$}{
        \For{$i \leftarrow 1$ \KwTo $n^3$}{
          Faz algo (custo $\Theta(1)$)\;
        }
        \textsc{Recursive}($n/3$)\;
      }
      \caption{\textsc{Recursive}($n$)}
    \end{algorithm}
    \begin{solution}
    Escreva a solução aqui.
  \end{solution}

  \end{parts}
  %28
  \question Assinale Verdadeiro (V) ou Falso (F). \textbf{Justifique}
  \begin{parts}
    \part \mbox{\begin{oneparcheckboxes}
        \correctchoice V \choice F
    \end{oneparcheckboxes}}
    O limite assintótico inferior para algoritmos de ordenação
    baseados em comparação é $\Omega(n \lg n)$. Um algoritmo de
    ordenação por comparação que faz $2T(n/2) + \Theta(1)$
    comparações no pior caso com certeza não efetua corretamente
    a ordenação para algumas instâncias.
    
    \part \mbox{\begin{oneparcheckboxes}
      \choice V \correctchoice F
    \end{oneparcheckboxes}}
    Suponha que iremos gerar $n$ números aleatórios no intervalo
    $[0 \ldots n^2]$. Considerando base decimal (isto é, um número $x$
    possui $\lfloor \log_{10} x \rfloor + 1$ dígitos na base $10$), é
    correto afirmar que o algoritmo \textsc{Radix Sort}
    faz a ordenação destes $n$ números em tempo $O(n)$.

    \part \mbox{\begin{oneparcheckboxes}
        \choice V \correctchoice F
    \end{oneparcheckboxes}}
    Visto que o limite assintótico inferior para algoritmos de ordenação
    baseados em comparação é $\Omega(n \lg n)$, não seria possível o
    desenvolvimento de um algoritmo de ordenação correto com complexidade
    de tempo $O(n \sqrt{n})$ no pior caso.
    
    \part \mbox{\begin{oneparcheckboxes}
        \choice V \correctchoice F
    \end{oneparcheckboxes}}
    Suponha que iremos gerar $n$ números aleatórios no intervalo
    $[0 \ldots n^2]$. É correto afirmar que o algoritmo \textsc{Counting Sort}
    faz a ordenação destes $n$ números em tempo $O(n)$.
    
    \part \mbox{\begin{oneparcheckboxes}
        \correctchoice V \choice F
    \end{oneparcheckboxes}}
    Suponha que iremos gerar $n$ números aleatórios no intervalo
    $[0 \ldots n^2]$. É correto afirmar que o algoritmo \textsc{Mergesort}
    faz a ordenação destes $n$ números em tempo $O(n \lg n)$.

    \part \mbox{\begin{oneparcheckboxes}
        \choice V \correctchoice F
    \end{oneparcheckboxes}}
    Não é possível construir um heap máximo com $n$ elementos em tempo
    $O(n)$. Pois para inserir um elemento no heap temos custo $O(\lg n)$ e,
    como temos $n$ elementos a serem inseridos, o custo total seria pelo
    menos $O(n \lg n)$.

    \part \mbox{\begin{oneparcheckboxes}
      \correctchoice V \choice F
    \end{oneparcheckboxes}}
    Em um heap binário, metade dos elementos do vetor são
    folhas. Se aplicarmos o procedimento \textsc{MaxHeapfy}
    para cada elemento da metade até o primeiro, então 
    ao fim teremos um Heap Máximo.

    \part \mbox{\begin{oneparcheckboxes}
        \correctchoice V \choice F
    \end{oneparcheckboxes}}
    É correto afirmar que: no melhor caso, o algoritmo \textsc{Insertion Sort}
    é mais eficiente que os algoritmos \textsc{Mergesort} e \textsc{Heapsort}. 
  \end{parts}
  \begin{solution}
    Escreva a solução aqui.
  \end{solution}
  %29
  \question Use o modelo de árvore de decisão para representar as
  comparações efetuadas pelo algoritmo \textsc{Insertion-Sort} para
  uma instância de entrada com quatro elementos.
  \begin{solution}
    Escreva a solução aqui.
  \end{solution}
  %30
  \question Use o modelo de árvore de decisão para representar as
  comparações efetuadas pelo algoritmo \textsc{Mergesort} para uma instância
  de entrada com três elementos.
  \begin{solution}
    Escreva a solução aqui.
  \end{solution}
  %31
  \question Use o modelo de árvore de decisão para representar as
  comparações efetuadas pelo algoritmo \textsc{Quicksort} para uma
  instância de entrada com três elementos.
  \begin{solution}
    Escreva a solução aqui.
  \end{solution}
  %32
  \question Dado um vetor de inteiros distintos e ordenados em
  ordem crescente $A = \{a_1, \ldots, a_n\}$:
  \begin{parts}
    \part Descreva um algoritmo que determina se existe um
    índice $i$ tal que $a_i = i$ em tempo $O(\lg n)$. Por exemplo, em
    $\{-10, -3, 3, 5, 7\}$, $a_3 = 3$. Em $\{2,3,4,5,6,7\}$ não existe
    tal $i$. Argumente que seu algoritmo está correto.
    \begin{solution}
      Escreva a solução aqui.
    \end{solution}
    \part Explique por que sua solução leva tempo $O(\lg n)$.
    \begin{solution}
      Escreva a solução aqui.
    \end{solution}
  \end{parts}
  %33
  \question Descreva um algoritmo baseado no paradigma de Divisão
  e Conquista que encontra o mínimo de um conjunto de $n$ números.
  Assuma que os elementos estão em um vetor $A = [1\ldots n]$.
  Mostre que o algoritmo está correto e analise sua complexidade.
  \begin{solution}
    Escreva a solução aqui.
  \end{solution}
  %34
  \question Descreva um algoritmo baseado no paradigma de Divisão
  e Conquista que encontra o segundo maior elemento de um conjunto
  de $n$ números. Assuma que os elementos estão
  em um vetor $A = [1\ldots n]$. Mostre que o algoritmo está
  correto e analise sua complexidade.
  \begin{solution}
    Escreva a solução aqui.
  \end{solution}
  %35
  \question Descreva um algoritmo que faz uso do procedimento
  \textsc{Partition} para encontrar o $k$-ésimo menor elemento.
  Isto é, o algoritmo recebe como entrada um vetor $A[1 \ldots n]$ e um
  um valor $1 \leq k \leq n$ e devolve qual seria este $k$-ésimo
  menor elemento. Por exemplo, se $k = 1$ o algoritmo deveria
  devolver o mínimo do vetor; se $k=n$ o algoritmo deveria devolver
  o máximo; para um $k=3$ devolveria o terceiro menor elemento.
  Analise seu algoritmo no pior e no melhor caso.
  \begin{solution}
    Escreva a solução aqui.
  \end{solution}
  %36
  \question Assuma que você possui $k$ vetores ordenados, cada um
  com $n$ elementos, e você precisa combiná-los em um único vetor
  ordenado com $k n$ elementos.
  \begin{parts}
  \part Usando o procedimento \textsc{Merge}, faça a intercalação
  do primeiro vetor com o segundo, então intercale o terceiro, depois
  o quarto e assim por diante. Qual a complexidade de tempo deste
  algoritmo, em função de $k$ e $n$?
  \begin{solution}
    Escreva a solução aqui.
  \end{solution}
  \part Apresente uma solução mais eficiente para este problema, por
  meio da técnica de Divisão e Conquista. Qual a complexidade de
  tempo de sua solução, em função de $k$ e $n$?
  \begin{solution}
    Escreva a solução aqui.
  \end{solution}
  \end{parts}
  %37
  \question Dado um vetor de números inteiros $A[1 \ldots n]$, determine
  quais elementos do vetor são únicos. Apresente um algoritmo eficiente.
  Faça uma análise de complexidade.
  \begin{solution}
    Escreva a solução aqui.
  \end{solution}
  %38
  \question Descreva um algoritmo de tempo $\Theta(n \lg n)$ que, dado
  um conjunto $S$ de $n$ números inteiros e outro número $x$, determine
  se existe dois elementos em $S$ cuja soma é exatamente $x$.
  \begin{solution}
    Escreva a solução aqui.
  \end{solution}
  %39
  \question Problema da moeda falsa. Dado um conjunto de $n$ moedas,
  $n-1$ delas verdadeiras (com mesmo peso) e uma falsa (mais leve),
  descreva um algoritmo eficiente (com tempo $o(n)$) para encontrar
  a moeda falsa.
  \begin{solution}
    Escreva a solução aqui.
  \end{solution}
  
\end{questions}

\end{document}
